% answers_short.tex —— 2024 年·新高考Ⅱ卷 · 简短答案速查
% 由 tools/make_short_answers.py 自答案要点 MD 逐题答案行生成，勿手改；
% 选择题字母经三源交叉核对（逐题行 / 速查表 / 解析 \gkanswer 框）。

\begin{gknotice}
\textbf{使用说明}：本节只给最终答案便于快速核对，完整推导见解析版。答案为编写者独立推导，\textbf{非官方标准答案}。
\end{gknotice}

\gkpart{单选题}
\begin{center}\begin{tabular}{ccccccccc}
\toprule
题号 & 1 & 2 & 3 & 4 & 5 & 6 & 7 & 8 \\
\midrule
答案 & C & B & B & C & A & D & B & C \\
\bottomrule
\end{tabular}\end{center}

\gkpart{多选题}
\begin{center}\begin{tabular}{cccc}
\toprule
题号 & 9 & 10 & 11 \\
\midrule
答案 & BC & ABD & AD \\
\bottomrule
\end{tabular}\end{center}

\gkpart{填空题}
\begin{itemize}[leftmargin=2.6em]
  \item[\textbf{12.}] $95$
  \item[\textbf{13.}] $-\dfrac{2\sqrt2}3$
  \item[\textbf{14.}] $24$；$112$
\end{itemize}

\gkpart{解答题}
\begin{itemize}[leftmargin=2.6em]
  \item[\textbf{15.}] (1) $A=\dfrac\pi6$；(2) 周长 $=2+3\sqrt2+\sqrt6$
  \item[\textbf{16.}] (1) $y=(\mathrm e-1)x-1$；(2) $a\in(1,+\infty)$
  \item[\textbf{17.}] (1) 见解析版；(2) $\dfrac{8\sqrt{65}}{65}$
  \item[\textbf{18.}] (1) $0.686$；(2) 见解析版
  \item[\textbf{19.}] (1) $x_2=3$，$y_2=0$；(2) 见解析版；(3) 见解析版
\end{itemize}
